\begin{resumo}
O crescimento acelerado da Creator Economy e a consolidação do ensino a distância 
no Brasil criaram uma demanda por ferramentas de gestão de conteúdo educacional 
que ferramentas genéricas como Trello, Notion e Asana não são capazes de suprir, 
por carecerem de estrutura semântica nativa para o domínio educacional. O presente 
trabalho descreve o desenvolvimento do Creator Workbench, uma plataforma web de 
Content Operations (ContentOps) projetada especificamente para centralizar e 
otimizar o pipeline de produção de criadores de conteúdo educacional digital. A 
metodologia adotada é o Scrum, com Sprints quinzenais, equipe de cinco integrantes 
com papéis técnicos distintos e repositório versionado no GitHub. A arquitetura do 
sistema segue o padrão Client-Server desacoplado, com front-end em React.js, 
back-end em Java com Spring Boot, banco de dados relacional MySQL normalizado até 
a Terceira Forma Normal e infraestrutura provisionada na Amazon Web Services. A 
plataforma oferece módulos de gestão de ideias, workflow Kanban com estados de 
produção, editor de roteiros, upload de ativos digitais integrado ao AWS S3, 
automações via scripts Python e estruturação curricular em Cursos, Módulos e Aulas. 
A segurança é implementada por autenticação JWT, criptografia BCrypt e comunicação 
HTTPS, em conformidade com a Lei Geral de Proteção de Dados Pessoais. A análise 
de viabilidade financeira demonstra sustentabilidade do modelo SaaS com três planos 
de assinatura, apresentando margem operacional superior a 90\% a partir do sétimo 
mês no cenário realista. Os resultados indicam que uma ferramenta com semântica 
nativa para o domínio educacional reduz a fragmentação operacional do criador e 
cria condições técnicas para a aplicação dos princípios andragógicos na produção 
de cursos digitais.

\vspace{\onelineskip}
\noindent
\textbf{Palavras-chave}: Creator Economy. Content Operations. Gestão de Conteúdo Educacional. Educação a Distância. SaaS.
\end{resumo}